(i) The one-loop coefficients are Fact \claimref{standard-model.imported-coefficient-sum-indices} applied to the component listings of Proposition \claimref{standard-model.table-exact} (gauginos one Weyl adjoint per factor, higgsinos a vector-like pair of doublets, the second doublet a complex scalar, the sfermions three scalar generations): the light sets add $\Dl = (\tfrac25, 2, 2)$, $(\tfrac25, \tfrac23, 0)$, $(\tfrac12, \tfrac56, 0)$ and $(\tfrac25, 2, 0)$ to $(\tfrac{41}{10}, -\tfrac{19}{6}, -7)$, and light plus heavy is the full content $(\tfrac52, \tfrac{25}{6}, 4)$ in every split. The two-loop matrices are the generic formula of Fact \claimref{standard-model.imported-two-loop-gauge-coefficients} on the listed Weyl fermions and complex scalars, which excludes gauge Yukawa couplings; the script rebuilds the Fact's Standard Model and supersymmetric matrices from components as controls, and prints the four middle-regime matrices. At $f = f'$ the middle regime has zero length, so the run is the one-threshold run of the full content; in Theorem \claimref{standard-model.two-threshold-criterion} the line $f(\Dl_{12} - R\Dl_{23}) + f'(\Dl'_{12} - R\Dl'_{23}) = Rb_{23} - b_{12}$ meets the diagonal at $f = (Rb_{23} - b_{12})/((\Dl + \Dl')_{12} - R(\Dl + \Dl')_{23})$, the same exact rational for all four splits, and at the inputs of Proposition \claimref{standard-model.spectrum-run} it is $1.0116\ldots > 1$, the threshold fraction of the full content in Proposition \claimref{standard-model.table-exact}(c). The script evaluates the four lines exactly, recovers the endpoints $(1, 0.99115)$ and $(0.4377, 0)$ of Proposition \claimref{standard-model.four-supersymmetric-splits}(d) and confirms that the lines of (a), (b) and (c) meet neither $f = 1$ nor $f' = 0$ inside the triangle; each passes through the diagonal at $1.0116$ and leaves the triangle on both sides.

(ii) A pair $(M_{I}, M_{I}')$ with $M_{I} = M_{I}'$ is admissible for a split exactly when the full content at the single threshold $M_{I}$ unifies, since the run is then identical; Proposition \claimref{standard-model.supersymmetric-row-two-loops} finds that it does at $M_{I} = 2.3$ TeV with $M_{U} = 1.1 \times 10^{16}$ GeV and $\alpha_{U}^{-1} = 25.7$, and states the spread $1.7$ to $3.0$ TeV across $\alpha_{s} = 0.1180 \pm 0.0009$. Since $M_{Z} < M_{I} = M_{I}' < M_{U}$, the point lies in the triangle $0 \le f' \le f \le 1$ for every split, so every split has a nonempty admissible set at two loops. Proposition \claimref{standard-model.proton-condition-candidates}(3) records that this point meets the proton condition.

(iii) to (v) are computed numerically, at record level, and are not part of the exact check; they are reproducible from the description that follows. The two-loop equations of Fact \claimref{standard-model.imported-two-loop-gauge-coefficients}, $da_{i}/d\ln\mu = -b_{i}/2\pi - \sum_{j}B_{ij}/(8\pi^{2}a_{j}) + c_{i}y_{t}^{2}/(32\pi^{3})$ with the middle-regime matrices of (i), are integrated from $M_{Z} = 91.1876$ GeV with the record's central values $\alpha_{\mathrm{em}}^{-1} = 127.951$, $\sin^{2}\theta_{W} = 0.23122$, $\alpha_{s} = 0.1180$, the top Yukawa running at one loop as the Fact gives it from $y_{t}(M_{Z}) = 0.95$, the thresholds as step functions and the matching $y_{t} \to y_{t}/\sin\beta$, $\sin\beta = 1$, at $M_{I}'$, by fourth-order Runge-Kutta in $\ln\mu$ with step $0.04$ (halving the step changes no printed digit). For each light threshold the heavy threshold is found by bisection on the mismatch $a_{2} - a_{3}$ at the crossing of $a_{1}$ and $a_{2}$. Controls: with the two-loop terms off and the rounded inputs of Proposition \claimref{standard-model.spectrum-run}, the integrator returns the exact one-loop endpoint $f' = 0.99115$ of split (d) at $f = 1$ that the check computes in rationals; with them on, the diagonal point is $M_{I} = M_{I}' = 2259$ GeV, $M_{U} = 1.125 \times 10^{16}$ GeV, $\alpha_{U}^{-1} = 25.736$, $f = 0.9011$, the stored row's values. The numbers obtained, without the gaugino Yukawas: far ends at $f = 1$ with the heavy set at $9.74 \times 10^{7}$, $7.92 \times 10^{5}$ and $1.60 \times 10^{9}$ GeV and $M_{U} = 2.67 \times 10^{16}$, $2.80 \times 10^{15}$, $4.76 \times 10^{14}$ GeV for (a), (b), (c); the $f' = 0$ end of (d) at a light set of $3.50 \times 10^{9}$ GeV with $M_{U} = 1.18 \times 10^{14}$ GeV and $\alpha_{U}^{-1} = 39.97$. A light threshold half an e-fold beyond the diagonal point, above it for (a) to (c) and below it for (d), admits no heavy threshold, which fixes the side of the diagonal on which each curve lies; the curves were traced by a scan of light thresholds at half-e-fold spacing, each admitting one heavy threshold. The proton figure of Fact \claimref{standard-model.imported-proton-condition} changes sign between light thresholds $450$ and $530$ GeV for (b) (heavy set $4.3$ to $3.2 \times 10^{4}$ GeV), between $1.1$ and $1.3$ TeV for (c) (heavy set $4.6$ to $2.3 \times 10^{4}$ GeV), and between $1.6$ and $1.8 \times 10^{4}$ GeV for (d) (heavy set $6.7$ to $8.2 \times 10^{4}$ GeV, $M_{U}$ from $5.9$ to $5.7 \times 10^{15}$ GeV); bisection locates the crossings stated in (iv). It is positive at both ends of (a) and at the diagonal point, and negative at the far ends of (b), (c) and (d). With the gaugino-Higgs-higgsino Yukawas included at tree-level matching the far end of (a) moves to $5.4 \times 10^{7}$ GeV and that of (d) to $3.39 \times 10^{9}$ GeV. At $\alpha_{s} = 0.1171$ and $0.1189$ the diagonal point is at $2990$ and $1714$ GeV and (d)'s far end at $3.87$ and $3.16 \times 10^{9}$ GeV. Two independent reproductions agree: integrating the two-loop terms along the one-loop path (leading logarithms) gives the diagonal point at $2.2$ TeV and the far ends within twenty percent in the heavy scale, and an adaptive integration (relative tolerance $10^{-11}$) reproduced every number above to the printed precision. The run's content, three derived generations with their lepton doublets, rests on Hypothesis (S) through Proposition \claimref{standard-model.spectrum-run}.
